\section{Accessing the fields}\label{sec:07-groups:05-accessing-fields:1}

\texttt{intGroup : Group Int} is an ordinary term, so its fields project out exactly as in Chapter 3.

\begin{lstlisting}
#eval intGroup.op 3 4        -- 7
#eval intGroup.id             -- 0
#eval intGroup.inv 5          -- -5

#check intGroup.assoc         -- a proof, for every a b c, of associativity
\end{lstlisting}

\begin{mathreading}

\texttt{intGroup.op} is \(\cdot\) (\texttt{intGroup.op 3 4} is \(3 + 4 = 7\)), \texttt{intGroup.id} is \(e = 0\), \texttt{intGroup.inv} is \((-)^{-1} =
-(-)\). \texttt{intGroup.assoc} projects a \emph{proof}: the element of \(\forall
a,b,c,\ (a\cdot b)\cdot c = a\cdot(b\cdot c)\) supplied when the group was built. Data-fields and proof-fields are accessed the same way, since a \texttt{structure} is, underneath, a dependent pair and both are coordinates of the same tuple.

\end{mathreading}

\textbf{Mathlib equivalent.} No field access to write at all. Once \texttt{Int} is an \href{https://loogle.lean-lang.org/?q=AddCommGroup}{\texttt{AddCommGroup}}, the ordinary \texttt{+}/\texttt{0}/\texttt{-\allowbreak{}} notations already resolve to the operations of that instance.

\begin{lstlisting}
#eval (3 : Int) + 4
#eval (0 : Int)
#eval -(5 : Int)
#check (add_assoc : (*$\forall$*) a b c : Int, (a + b) + c = a + (b + c))
\end{lstlisting}

Same contrast as Section 3: \texttt{intGroup.op}/\texttt{.id}/\texttt{.inv} are projections out of a hand-built bundle; \texttt{+}/\texttt{0}/\texttt{-\allowbreak{}} are notation the typeclass system has already wired to the right instance, with the underlying ``which \texttt{AddCommGroup} instance'' bookkeeping invisible unless sought (for instance, with \texttt{\#print}).
